\documentclass[12pt, a4paper]{article}
\usepackage[utf8]{inputenc}
\usepackage[brazil]{babel}
\usepackage[margin=2.5cm]{geometry}
\usepackage{hyperref}
\usepackage{listings}
\usepackage{xcolor}
\usepackage{graphicx}

% Configuração para exibição de código C e Makefile
\lstset{
    basicstyle=\ttfamily\small,
    keywordstyle=\color{blue}\bfseries,
    stringstyle=\color{red},
    commentstyle=\color{green!60!black},
    numbers=left,
    numberstyle=\tiny\color{gray},
    stepnumber=1,
    frame=single,
    breaklines=true,
    captionpos=b
}

\title{\textbf{Relatório Técnico: Sistema de Processamento de Áudio (\texttt{audularia})}}
\author{TODO\\TODO\\TODO\\TODO}}
\date{\today}

\begin{document}

\maketitle

\tableofcontents
\newpage

\section{Descrição do Problema e do Domínio Escolhido}
O projeto consiste em uma biblioteca em linguagem C desenvolvida para o processamento, conversão e manipulação de arquivos de áudio. O objetivo principal é permitir carregar, alterar suas propriedades  e aplicar filtragens no domínio da frequência utilizando a Transformada Rápida de Fourier (FFT).

\section{Arquitetura do Software}
A arquitetura do projeto é dividida em módulos independentes organizados em pastas:
\begin{itemize}
    \item \texttt{include/}: Contém os arquivos de cabeçalho (\texttt{.h}) que definem as interfaces públicas.
    \item \texttt{src/}: Contém as implementações (\texttt{.c}) do núcleo da biblioteca.
    \item \texttt{app/}: Contém a aplicação cliente (\texttt{main.c}).
    \item \texttt{test/}: Contém os testes automatizados para validação das funções.
    \item \texttt{lib/}: Armazena as bibliotecas geradas (\texttt{.a} estática e \texttt{.so} dinâmica).
    \item \texttt{bin/}: Armazena os executáveis finais.
\end{itemize}

\section{Organização e Responsabilidade dos Módulos}
Os módulos estão divididos da seguinte forma:
\begin{enumerate}
    \item \textbf{\texttt{audio} (\texttt{audio.h} / \texttt{audio.c}):} Gerencia a estrutura do áudio, alocação de memória e conversão de formatos (WAV, MP3 e RAW).
    \item \textbf{\texttt{operacoes} (\texttt{operacoes.h} / \texttt{operacoes.c}):} Implementa edições diretas nas amostras, como mudança de intensidade (volume), velocidade, inversão, corte e mixagem de dois áudios.
    \item \textbf{\texttt{fft} (\texttt{fft.h} / \texttt{fft.c}):} TODO
    \item \textbf{\texttt{fft\_audio} (\texttt{fft\_audio.h} / \texttt{fft\_audio.c}):} TODO
\end{enumerate}

\section{Descrição da API Pública}
TODO


\section{Descrição da Aplicação Cliente}
A aplicação cliente fica no arquivo \texttt{app/main.c}. Ela é responsável por interagir diretamente com o usuário/sistema, chamando a API pública para carregar um arquivo (como o \texttt{Super\_mario.raw}), aplicar os efeitos ou filtros desejados e salvar/executar o resultado final.

\section{Processo de Compilação Utilizando Makefile}
O projeto utiliza um arquivo \texttt{Makefile} automatizado para compilação. As variáveis principais do \texttt{Makefile} incluem:
\begin{itemize}
    \item \texttt{CC = gcc}: Compilador C.
    \item \texttt{CFLAGS = -Wall -Wextra -Iinclude -fPIC}: Flags para exibir alertas, incluir o diretório de cabeçalhos e permitir código de posição independente (necessário para bibliotecas dinâmicas).
\end{itemize}
As regras automatizam a criação das pastas \texttt{build/}, \texttt{bin/} e \texttt{lib/} antes de compilar os arquivos objeto.

\section{Estratégia de Compilação Incremental}
TODO

\section{Geração e Utilização da Biblioteca Estática}
A biblioteca estática é gerada no diretório \texttt{lib/libaudio.a} usando o utilitário \texttt{ar}:
\begin{lstlisting}[language=make]
$(STATIC_LIB): $(OBJS)
    @mkdir -p $(LIB_DIR)
    $(AR) rcs $@ $^
\end{lstlisting}
Ao compilar o programa cliente com a biblioteca estática, todo o código necessário da biblioteca é "copiado" diretamente para dentro do executável final.

\section{Geração e Utilização da Biblioteca Dinâmica}
A biblioteca dinâmica é gerada no arquivo \texttt{lib/libaudio.so} usando a flag \texttt{-shared}:
\begin{lstlisting}[language=make]
$(DYNAMIC_LIB): $(OBJS)
    @mkdir -p $(LIB_DIR)
    $(CC) -shared $^ -o $@
\end{lstlisting}
Para rodar a aplicação vinculada dinamicamente, o executável apenas referencia a biblioteca em tempo de execução, reduzindo o tamanho do arquivo final.

\section{Comparação Entre Ligação Estática e Dinâmica}
\begin{itemize}
    \item \textbf{Ligação Estática (\texttt{.a}):} 
    \begin{itemize}
        \item \emph{Vantagem:} O executável é autossuficiente (não depende de arquivos externos ao rodar).
        \item \emph{Desvantagem:} Gera arquivos binários maiores e requer recompilar o cliente se a biblioteca mudar.
    \end{itemize}
    \item \textbf{Ligação Dinâmica (\texttt{.so}):} 
    \begin{itemize}
        \item \emph{Vantagem:} Binários menores e possibilidade de atualizar a biblioteca sem recompilar a aplicação cliente.
        \item \emph{Desvantagem:} O executável exige que a biblioteca esteja instalada no sistema no momento da execução.
    \end{itemize}
\end{itemize}

\section{Estratégia e Resultados dos Testes}
Os testes foram organizados em \texttt{test/test\_run.c}. A estratégia adotada baseou-se em testes unitários para validar individualmente:
\begin{enumerate}
    \item Criação e destruição de estruturas de áudio sem vazamento de memória (\emph{memory leaks}).
    \item Validação matemática de filtros e da transformada de Fourier.
    \item Verificação de condições de erro em entradas inválidas (ex: frequências negativas ou ponteiros nulos).
\end{enumerate}

\section{Exemplo de Teste que Inicialmente Falhou e Posterior Correção}
TODO

\section{Estratégia de Versionamento Utilizando Git}
TODO

\section{Descrição da Branch, Merge e Tag v1.0}
TODO

\section{Dificuldades Encontradas e Decisões de Projeto}
TODO

\section{Resultados Obtidos e Possíveis Melhorias}
O sistema final se mostrou estável e capaz de processar e manipular arquivos sonoros de forma rápida e modular. Como melhorias futuras, destacam-se:
\begin{itemize}
    \item Suporte nativo a múltiplos canais (áudio estéreo).
    \item Implementação de interface gráfica para visualização das formas de onda e espectro de frequência.
\end{itemize}

\end{document}